%! Systems Involving Quadratics — Unit 3, Lesson 3.6 — Warm-Up (Answer Key)
\documentclass[10pt]{article}
\usepackage{algebra2-article}
\usepackage{algebra2-key}   % pulls in -boxes; do NOT also load -boxes
\newcommand{\TallMath}[1]{$\displaystyle #1\rule[-1.4em]{0pt}{3.2em}$}

\begin{document}

\pageheader{Unit 3, Lesson 3.6}{Warm-Up --- Answer Key}
\namedateperiod

\vspace{0.10in}

\begin{notesbox}{Getting Ready: Skills We'll Use Today}
\small Every system today collapses to \emph{one} quadratic equation. These three skills are the tools that
finish it --- and the habit that confirms an answer.

\vspace{4pt}
\textbf{1. Solve by factoring} (Lesson 3.2). Set each factor to zero.
\[
  x^2-3x=0:\ \ x=\ans{$0,\ 3$} \qquad\qquad x^2-2x-8=0:\ \ x=\ans{$4,\ -2$}
\]

\vspace{4pt}
\textbf{2. Count the roots with the discriminant} (Lesson 3.5). Compute $b^2-4ac$; state \emph{two real},
\emph{one repeated}, or \emph{two complex} --- \emph{no solving}.
\[
  x^2-2x+1=0:\ \ b^2-4ac=\ans{$0$}\ \Rightarrow\ \ans{one repeated real}
\]
\[
  x^2-2x+3=0:\ \ b^2-4ac=\ans{$-8$}\ \Rightarrow\ \ans{two complex}
\]

\vspace{4pt}
\textbf{3. Does the point work?} Check whether $(3,0)$ makes \emph{each} equation true (write yes/no).
\[
  y=x^2-2x-3:\ \ans{yes} \qquad\qquad y=x-3:\ \ans{yes}
\]
\emph{Today: a point is a solution of a \textbf{system} only when it makes \textbf{both} equations true at once.}
\end{notesbox}

\begin{teachernote}
Each item is a direct prerequisite for solving a system. Item~1 is the factoring that finishes a collapsed
system --- note $x^2-3x=0$ needs a GCF ($x(x-3)$) and has the easy-to-miss root $x=0$. Item~2 previews ``how many
intersection points'': the collapsed quadratic's discriminant is $0$ (line tangent to parabola, \emph{one} point)
and $-8<0$ (line misses, \emph{no} points) --- exactly the $y=-4$ and $y=-6$ hook cases. Item~3 builds the
verification habit: $(3,0)$ satisfies \emph{both} the anchor parabola ($9-6-3=0$) and the line ($3-3=0$), so it is
a solution of the system; stress that a point on only one graph does \emph{not} qualify. Debrief item~3 last.
\end{teachernote}

\end{document}
